\sgaStableTarget{sga6.e05.expose.occ2}\sgaStableTarget{sga6.e05.def.6-2}\textbf{Definition 6.2.} \uline{The functor defined in} 6.1 \uline{on the category of} \uline{graded} @@M0@@-\uline{algebras}, \uline{with values in the category of} @@M1@@-@@M2@@-\uline{algebras}, \uline{is called the} \uline{Chern ring}; \uline{its value on} @@M3@@ \uline{is denoted by} @@M4@@ (\uline{or}, \uline{if} \uline{there is a risk of confusion}, @@M5@@).

From formulas (6.1.1), (6.1.2), (6.1.4), and (3.6.2), we deduce the relations:
\sgaStableTarget{sga6.e05.eq.6-2-1}@@M6@@
relations which served to define the Chern ring in [1]. Conversely, from these formulas one can recover relations (6.1.1) and (6.1.4).

Let us indicate why the Chern ring is introduced. We place ourselves in the situation of Example 3.9, and denote by @@M7@@ the associated graded ring (for the filtration defined in 3.10) of the @@M8@@ of a commutatively locally ringed topos (which we suppose to be a @@M9@@-ring for simplicity); @@M10@@ is the augmentation ring. We shall see below that, for every element of an augmented @@M11@@-ring, one can define Chern classes with values in the associated graded ring (6.7). To every locally free sheaf of finite type, we can therefore associate an element of @@M12@@ whose component on the factor @@M13@@ is the rank of the sheaf, and whose component on the factor @@M14@@ is the sum of the Chern classes. Formulas (6.2.1) are then the usual formulas giving the Chern classes and the rank of the tensor product of two invertible sheaves and of the exterior powers of an invertible sheaf (that is, the map just defined from @@M15@@ to @@M16@@ is a @@M17@@-homomorphism).

